\section{Smooth approximation in the actual theta graph norm}
\label{sec:spectral-prolate-graph}

The following estimate concerns the original sharp prolate
projection and the closed arithmetic theta map.  For smooth
input it needs only the unconditional convexity bound, not
theorem~A.  It is also an estimate for families of columns,
which is useful before any arithmetic Gram normalization.

\begin{theorem}[AC-S29: quantitative theta graph approximation]
\label{AC-S29}
Fix a nonprincipal primitive Dirichlet character $\chi$,
with parity $\epsilon$.  Use the Fourier convention
$\widehat f(\xi)=\int_{\mathbb R}f(u)e^{-2\pi i u\xi}\,du$.
On the raw parity space, with norm
$\|f\|_{\mathcal H}^2=\int_0^\infty|f(u)|^2du$, put
\[
 P=1_{[-1,1]},\qquad
 Q_c=\mathcal F^{-1}1_{[-c,c]}\mathcal F,
             \qquad S_c=PQ_cP.                     \tag{SG.1}
\]
Let $\Pi_c$ retain an initial segment of the prolate
eigenvectors in this parity, and suppose that, uniformly
for the values of $c\geq1$ under consideration,
\[
 \min_{\rm retained}\nu_j\geq\eta,
       \qquad\max_{\rm discarded}\nu_j\leq1-\eta
       \quad\text{for some }\eta>0.                 \tag{SG.2}
\]
The original full cutoff satisfies these inequalities for
all sufficiently large $c$, by the concentration bounds
already used in AC-T29 and AC-T33.

Define the dilation Sobolev norm by
\[
 (\mathcal Uf)(t)=e^{t/2}f(e^t),\qquad
       \|f\|_{H_D^s}=\|\mathcal Uf\|_{H^s(\mathbb R)}.
                                                               \tag{SG.3}
\]
The closed theta operator $\mathcal E_\chi$ is the
Mellin multiplier of AC-T21 and AC-T25, agreeing on
its smooth core with
$E_\chi f(x)=\sqrt x\sum_{n\geq1}\chi(n)f(nx)$.
For $u=3/8$, $\vartheta=5/16$ and $f=Pf\in H_D^u$,
\[
 \begin{split}
 \|\mathcal E_\chi(1-\Pi_c)f\|_2
 \leq{}&C_{\chi,\eta}
        \|(1-Q_c)f\|_{\mathcal H}^{1/6}\\
 &\quad\cdot
   \bigl(\|f\|_{H_D^u}+(1+c)^u\|f\|_{\mathcal H}\bigr)^{5/6}.
 \end{split}                                        \tag{SG.4}
\]
All vectors on the left belong to the closed theta domain.

There is a corresponding column estimate.  If
$T:\mathcal A\to\mathcal H$ is Hilbert--Schmidt,
$PT=T$, and $T$ is Hilbert--Schmidt into $H_D^u$, then
\[
 \begin{split}
 \|\mathcal E_\chi(1-\Pi_c)T\|_2
 \leq{}&C_{\chi,\eta}
              \|(1-Q_c)T\|_2^{1/6}\\
 &\quad\cdot
 \bigl(\|T\|_{2,\mathcal A\to H_D^u}
                     +(1+c)^u\|T\|_2\bigr)^{5/6}.
 \end{split}                                        \tag{SG.5}
\]
Here the subscript $2$ on operators denotes the
Hilbert--Schmidt norm, while the left side of (SG.4)
is the target $L^2(dx/x)$ norm.
In particular, if a family $T_c$ satisfies for every $m$
\[
 \|(1-Q_c)T_c\|_2\leq C_m c^{-m}A_c,
 \qquad
 \|T_c\|_{2,\mathcal A\to H_D^u}
                 +c^u\|T_c\|_2\leq Cc^u A_c,
                                                               \tag{SG.6}
\]
then $\|\mathcal E_\chi(1-\Pi_c)T_c\|_2
\leq C_Nc^{-N}A_c$ for every prescribed $N$.
The constants in (SG.6) must be uniform for the family;
no such uniformity is inferred from pointwise smoothness.

For every integer $m\geq1$, let $X_m$ be the parity
subspace of $H^m(\mathbb R)$ supported in $[-1,1]$,
with its ordinary spatial Sobolev norm.  Then
\[
 \|\mathcal E_\chi(1-\Pi_c)\|_{X_m\to L^2(dx/x)}
                \leq C_{\chi,\eta,m}
                         c^{-m/6+5/16}.             \tag{SG.7}
\]
Consequently a fixed $f\in C_c^\infty((-1,1))$ of the
prescribed parity is approximated by $\Pi_cf$ in
both raw $L^2$ norm and theta graph norm faster than
any prescribed negative power of $c$.
If $T:\mathcal A\to X_m$ belongs to a Schatten class
$\mathcal S_p$, $1\leq p\leq\infty$, composition in
(SG.7) gives the same bound multiplied by
$\|T\|_{p,\mathcal A\to X_m}$.
Thus fixed smooth compact kernel operators whose output
has the prescribed parity and is supported strictly inside
$(-1,1)$ have the same arbitrary-power convergence in
trace norm and Hilbert--Schmidt norm.
\end{theorem}
\begin{proof}
\proofstep{1. Raw concentration and the sharp endpoint.}
Expand $f=Pf$ in the time-supported prolate basis.
For every discarded index, $1-\nu_j\geq\eta$.
The spectral expansion and the projection identity give
\[
 \begin{split}
 \|(1-\Pi_c)f\|_{\mathcal H}^2
 &\leq\eta^{-1}\langle f,(1-S_c)f\rangle\\
 &=\eta^{-1}\|(1-Q_c)f\|_{\mathcal H}^2.
 \end{split}                                        \tag{SG.8}
\]
The same inequality holds after summing over an
orthonormal basis of a Hilbert--Schmidt input space.

For $\psi\in\operatorname{ran}\Pi_c$, put
$\phi=S_c^{-1}\psi$ on the retained finite-dimensional
space.  Then $\|\phi\|\leq\eta^{-1}\|\psi\|$ and
$\psi=Q_c\phi$ inside $(-1,1)$.
Band limitation and Plancherel therefore imply
\[
       \|\psi'\|_{L^2(0,1)}
             \leq2\pi c\eta^{-1}\|\psi\|_{\mathcal H}.
                                                               \tag{SG.9}
\]
The positive-half-line normalization follows from parity.
Set $v(t)=e^{t/2}\psi(e^t)$ for $t<0$.
Then $\|v\|_2=\|\psi\|$ and
$\|v'\|_2\leq C_\eta(1+c)\|\psi\|$.
The zero extension of $v$ above zero has a possible jump.
For $0<s<1/2$, the explicit zero-extension estimate
AC-T28, equation (TB.33), applies with this endpoint:
\[
 \|\psi\|_{H_D^s}
   \leq C_{\eta,s}(1+c)^s\|\psi\|_{\mathcal H}.
                                                               \tag{SG.10}
\]
To recall its endpoint justification, integration by
parts gives $i\tau\widehat v$ as the transform of
$v'$ together with one term $v(0)$, and
$|v(0)|^2\leq2\|v\|_2\|v'\|_2$.
Splitting at $1+\|v'\|_2/\|v\|_2$ proves the estimate;
the boundary tail is integrable exactly when $s<1/2$.
Thus (SG.10) includes the actual sharp time cutoff.

\proofstep{2. The theta graph estimate.}
The unconditional bound already proved in AC-T25,
equation (TB.16), is
\[
 |L(1/2+i\tau,\chi)|
                    \leq C_\chi(1+|\tau|)^{5/16}.
                                                               \tag{SG.11}
\]
It follows from the absolutely convergent line
$\operatorname{Re}s=17/16$, the functional equation
on $\operatorname{Re}s=-1/16$, and strip convexity;
it uses no zero-location assumption.
By the closed Mellin realization,
\[
       \|\mathcal E_\chi h\|_2
                      \leq C_\chi\|h\|_{H_D^{\vartheta}},
                  \qquad\vartheta=5/16.             \tag{SG.12}
\]
In particular every vector in that Sobolev space
belongs to the closed multiplier domain.

Let $h=(1-\Pi_c)f$ and $u=3/8$.
Fourier H\"older interpolation gives
\[
 \|h\|_{H_D^{\vartheta}}
       \leq\|h\|_{\mathcal H}^{1-\vartheta/u}
                           \|h\|_{H_D^u}^{\vartheta/u}.
                                                               \tag{SG.13}
\]
Here $\vartheta/u=5/6$.
Equations (SG.8) and (SG.10), and
$\|\Pi_cf\|\leq\|f\|$, prove (SG.4).
For (SG.5), perform exactly the same H\"older inequality
on the product of Fourier measure and counting measure
over an orthonormal basis of $\mathcal A$.
This gives
\[
 \|h\|_{2,\mathcal A\to H_D^{\vartheta}}
 \leq\|h\|_2^{1-\vartheta/u}
                     \|h\|_{2,\mathcal A\to H_D^u}^{\vartheta/u}
\]
for a column operator $h$.
The summed forms of (SG.8) and (SG.10) prove the
stated bound.  Finite-column approximation extends it
to all the stated Hilbert--Schmidt maps, using
closedness of the theta multiplier.
Substitution of (SG.6) gives
$C_m c^{-m/6+5/16}A_c$; choosing $m$ sufficiently
large proves its consequence.

\proofstep{3. Spatial smoothness and Schatten composition.}
For $f\in X_m$, Plancherel bounds its Fourier tail by
\[
       \|(1-Q_c)f\|\leq C_m c^{-m}\|f\|_{H^m}.
\]
Also $\|f\|_{H_D^u}\leq C\|f\|_{H^1}$.
Indeed differentiating $e^{t/2}f(e^t)$ gives the
two terms $e^{t/2}f(e^t)/2$ and $e^{3t/2}f'(e^t)$;
the latter has squared integral
$\int_0^1u^2|f'(u)|^2du\leq\|f'\|_2^2$.
The trace at $u=1$ is zero for the full-line
$H^1$ zero extension, so no boundary distribution
appears in this assertion.
Consequently the second factor in (SG.4) is bounded
by $C_m c^u\|f\|_{H^m}$ for $c\geq1$.
This proves (SG.7), including its exponent
$-m/6+(5/6)(3/8)=-m/6+5/16$.
The raw norm estimate follows separately from (SG.8).

The Schatten conclusion is the operator-ideal inequality
$\|AT\|_p\leq\|A\|\|T\|_p$ applied to the bounded map
in (SG.7).
For a fixed smooth kernel compactly supported in both
variables, a Fourier expansion in its input variable
has coefficients rapidly decreasing in every output
$H^m$ norm, by integration by parts.
The resulting rank-one series is absolutely summable
as a map into $X_m$ whenever its output support is
contained in $(-1,1)$.
Thus it is trace class into every such $X_m$, proving
the last assertion.
\end{proof}

\paragraph{Exact physical scaling and remaining scope.}
For the actual time radius $\Lambda$ and frequency radius
$\Lambda/q$, one has $c=\Lambda^2/q$ after the unitary
dilation $D_\Lambda f(u)=\Lambda^{-1/2}f(u/\Lambda)$.
This dilation preserves all the $H_D^s$ norms, and
$\mathcal E_\chi D_\Lambda=V_\Lambda\mathcal E_\chi$
with $V_\Lambda F(x)=F(x/\Lambda)$ unitary on the target.
Therefore all estimates above apply to the actual
conductor-normalized finite prolate family, with that
exact substitution.  They control smooth profiles or
families satisfying the displayed norm hypotheses.
They do not assert uniform boundedness of $\Pi_c$ on
$H_D^s$, convergence of every varying boundary family,
or the needed comparison of arithmetic polar projections.
